% Abstract text — shared across all versions
% Do not add \begin{abstract} here — the wrapper handles that

The Voynich Manuscript (Beinecke MS~408, radiocarbon-dated 1404--1438~CE) is a 234-page illustrated codex written in an unidentified script of approximately 166 distinct glyphs.  We develop a four-stage computational pipeline---statistical fingerprinting, generative model comparison, syllabic segmentation, and Bayesian decipherment---applied to the STA1~2.0 corpus (37,087 words, 157,254 glyph tokens).  The fingerprint reveals a unique U-shaped entropy curve (bounce $= 0.813$, absent from all controls) and positional mutual information (0.658~bits) twice that of any tested language.  Comparison of eleven generative models identifies a syllabic-partition mechanism (SilPart, $d = 0.445$) as the only model reproducing the full fingerprint; four published hypotheses are quantitatively eliminated.  Bayesian ranking across six candidate languages places Hebrew and Aramaic as statistically indistinguishable top candidates (gap $\sim$0.01~nats, within bootstrap CI spanning zero).  However, nineteen decipherment models---spanning trigram and neural LMs, single and triple mapping tables, character and syllable targets, morphological role indexing, and three additional cross-linguistic targets (Tibbonid Hebrew, Judeo-Arabic, Classical Ge'ez)---all converge on the same structural ceiling: approximately 68\% $n$-gram realism and 3--5 root matches at best.  No model produces readable text.  The Ge'ez experiment ($|\Sigma|/|\mathcal{A}| \approx 0.88$, the first near-1:1 test) confirms that the barrier is not an artefact of alphabet-size mismatch.  Classical cipher-type diagnostics (IC, Kasiski, bigram entropy, transposition) further exclude polyalphabetic, polygraphic, and transposition ciphers, establishing a central paradox: the text behaves statistically like natural language but resists all nineteen monoalphabetic models.  We establish quantitative falsification criteria (bounce $> 0.8$, MI$_{\text{pos}} > 0.5$) for evaluating future proposals.  The convergence of nineteen independent strategies provides strong evidence that the structural ceiling is a property of the Voynich text itself---not of any particular modelling choice---and narrows the remaining viable hypotheses to substitution to an untested language or a constructed language.  A companion study analyses the encoding mechanism directly, finding a two-tier combinatorial structure consistent with an encyclopaedic code.
